% !PRINTED 468
\weberpage{468}

According to the formulae \S.\,115, (14) one obtains first of all very simply:

\[
\mu = \frac{2\varepsilon}{7},
\]

and for $\lambda$ one obtains, if one substitutes the values \S.\,115, (11) for $\alpha, \beta, \gamma$, and uses the formulae \S.\,115, (12), (13):

\weq{5}{\begin{aligned}
\lambda &= -\varepsilon\,\frac{1+\sqrt{-7}}{14} \\
\frac{\mu}{\lambda} &= -\frac{1-\sqrt{-7}}{2}.
\end{aligned}}

If therefore one sets

\[
A_1^2+A_2^2+A_3^2 = \lambda z,
\]

then it follows that

\weq{6}{z = x_1^2+x_2^2+x_3^2 - \frac{1-\sqrt{-7}}{2}(x_2x_3+x_3x_1+x_1x_2),}

and this function we shall introduce as a root of the resolvent of degree 7. The remaining roots are obtained from it by the substitutions $\tau^r$, so that they all, in the common form

\weq{7}{\begin{aligned}
z_r &= \varepsilon^{2r}x_1^2 + \varepsilon^{4r}x_2^2 + \varepsilon^{r}x_3^2 \\
&\quad - \frac{1-\sqrt{-7}}{2}(\varepsilon^{-r}x_2x_3+\varepsilon^{-2r}x_3x_1+\varepsilon^{-4r}x_1x_2) \\
&\qquad r=0, 1, 2, 3, 4, 5, 6
\end{aligned}}

are contained.

The coefficients of the equation of degree 7 whose roots are the seven quantities $z_r$ are invariants of our group, whose degree can easily be specified.

For if $a_\nu$ is the coefficient of $z^{7-\nu}$ in this equation, after the coefficient of the seventh power has been brought to 1, then $a_\nu$ is an invariant of degree $2\nu$. Hence first of all we must have $a_1=0$, because there is no quadratic invariant. The remaining coefficients are expressed linearly and homogeneously in terms of the following combinations of invariants:

\weq{8}{\begin{aligned}
a_2 &\text{ by } f, \\
a_3 &\text{ \textquotedbl} \ \varDelta, \\
a_4 &\text{ \textquotedbl} \ f^2, \\
a_5 &\text{ \textquotedbl} \ \varDelta f, \\
a_6 &\text{ \textquotedbl} \ \varDelta^2, f^3, \\
a_7 &\text{ \textquotedbl} \ C, \varDelta f^2
\end{aligned}}
